\documentclass[10pt,a4paper]{amsart}
\usepackage{amsmath,amssymb,amsthm,graphicx}
% Retained minimal preamble for the Stacks Project chapter varieties.tex.
% Verbatim excerpt of the chapter preamble: Xy-pic setup (native preamble lines 14-18), the native theorem
% declarations with the mathematical macros (native preamble lines 154-217), and the project own no-op
% definitions of the reference/slogan/history display environments (native wrapper.tex lines 1-3) that the
% counted statement uses for its slogan. Copied mechanically; nothing added, removed or redefined.
% For commutative diagrams we use Xy-pic
\usepackage[all]{xy}

% We use 2cell for 2-commutative diagrams.
\xyoption{2cell}
% Theorem environments.
%
\theoremstyle{plain}
\newtheorem{theorem}[subsection]{Theorem}
\newtheorem{proposition}[subsection]{Proposition}
\newtheorem{lemma}[subsection]{Lemma}

\theoremstyle{definition}
\newtheorem{definition}[subsection]{Definition}
\newtheorem{example}[subsection]{Example}
\newtheorem{exercise}[subsection]{Exercise}
\newtheorem{situation}[subsection]{Situation}

\theoremstyle{remark}
\newtheorem{remark}[subsection]{Remark}
\newtheorem{remarks}[subsection]{Remarks}

\numberwithin{equation}{subsection}

% Macros
%
\def\lim{\mathop{\mathrm{lim}}\nolimits}
\def\colim{\mathop{\mathrm{colim}}\nolimits}
\def\Spec{\mathop{\mathrm{Spec}}}
\def\Hom{\mathop{\mathrm{Hom}}\nolimits}
\def\Ext{\mathop{\mathrm{Ext}}\nolimits}
\def\SheafHom{\mathop{\mathcal{H}\!\mathit{om}}\nolimits}
\def\SheafExt{\mathop{\mathcal{E}\!\mathit{xt}}\nolimits}
\def\Sch{\mathit{Sch}}
\def\Mor{\mathop{\mathrm{Mor}}\nolimits}
\def\Ob{\mathop{\mathrm{Ob}}\nolimits}
\def\Sh{\mathop{\mathit{Sh}}\nolimits}
\def\NL{\mathop{N\!L}\nolimits}
\def\CH{\mathop{\mathrm{CH}}\nolimits}
\def\proetale{{pro\text{-}\acute{e}tale}}
\def\etale{{\acute{e}tale}}
\def\QCoh{\mathit{QCoh}}
\def\Ker{\mathop{\mathrm{Ker}}}
\def\Im{\mathop{\mathrm{Im}}}
\def\Coker{\mathop{\mathrm{Coker}}}
\def\Coim{\mathop{\mathrm{Coim}}}

% Boxtimes
%
\DeclareMathSymbol{\boxtimes}{\mathbin}{AMSa}{"02}

%
% Macros for moduli stacks/spaces
%
\def\QCohstack{\mathcal{QC}\!\mathit{oh}}
\def\Cohstack{\mathcal{C}\!\mathit{oh}}
\def\Spacesstack{\mathcal{S}\!\mathit{paces}}
\def\Quotfunctor{\mathrm{Quot}}
\def\Hilbfunctor{\mathrm{Hilb}}
\def\Curvesstack{\mathcal{C}\!\mathit{urves}}
\def\Polarizedstack{\mathcal{P}\!\mathit{olarized}}
\def\Complexesstack{\mathcal{C}\!\mathit{omplexes}}
% \Pic is the operator that assigns to X its picard group, usage \Pic(X)
% \Picardstack_{X/B} denotes the Picard stack of X over B
% \Picardfunctor_{X/B} denotes the Picard functor of X over B
\def\Pic{\mathop{\mathrm{Pic}}\nolimits}
\def\Picardstack{\mathcal{P}\!\mathit{ic}}
\def\Picardfunctor{\mathrm{Pic}}
\def\Deformationcategory{\mathcal{D}\!\mathit{ef}}
\newenvironment{reference}{}{}
\newenvironment{slogan}{}{}
\newenvironment{history}{}{}
\usepackage{cleveref}
\crefname{originallemma}{original lemma}{original lemmas}
\Crefname{originallemma}{Original lemma}{Original lemmas}
\makeatletter
\newlabel{limits-lemma-descend-opens}{{01Z4}{source}{}{sourcepointer-1}{}}
\newlabel{limits-lemma-descend-opens@cref}{{[originallemma][0][]01Z4}{[1][0][]source}{}{}{}}
\makeatother
\title{A quasi-compact open of a base change descends to a finitely generated subalgebra}
\author{The Stacks Project Authors}
\date{stacks-project:varieties.tex:local-snapshot-20261011; GNU FDL 1.2 or later}
\begin{document}
\maketitle
\noindent\textit{Editorial note.} All delivered mathematical text is the author's own, copied from the pinned Stacks Project chapter file varieties.tex; no statement or proof step was rewritten, reordered or removed.\par\medskip
\noindent\textit{Editorial note.} The proof opens by remarking that the quasi-separated case is a known descent lemma; the remark is reproduced verbatim and the case actually proved here is the general one, with the argument written out inside the counted span.\par\medskip
\noindent\textit{Editorial note.} The counted spans carry 1 navigation stub and no bibliography citation.\par\medskip
\noindent\textit{Editorial note.} The preamble is the chapter's retained minimal preamble (sources/stacks-varieties/retained-preamble.tex lines 1-77), a verbatim excerpt of the native preamble.\par\medskip
\noindent\textit{Editorial note.} The label lemma-descend-open is defined only in varieties.tex, so no cross-chapter restatement had to be weighed for this card.\par\medskip
\noindent\textit{Allowed standard prerequisites (editorial).}\begin{itemize}
\item Descent of quasi-compact opens through a filtered colimit of rings. Limits, Lemma limits-lemma-descend-opens, invoked at proof lines 823-824. Supplies the result in the quasi-separated case and the model for the argument given here.
\item Standard opens in an affine scheme and the compatibility of localization with base change. Standard, used at proof lines 826-841. Supplies the description of the open V as a union of standard opens and the pointwise comparison of fibres.
\end{itemize}\medskip
\ifdefined\hypertarget\hypertarget{sourcepointer-1}{}\fi
\noindent\textit{Source reference.} Original lemma 01Z4: sources/stacks-chow/tags.txt, source lines 2579--2579. Pointer outside the counted span, invoked once for context; the load-bearing argument of the counted statement is the author's proof reproduced in full in the proof span.\par\medskip
\begin{lemma}
\label{lemma-descend-open}
Let $k$ be a field.
Let $X$ be a scheme over $k$.
Let $A$ be a $k$-algebra.
Let $V \subset X_A$ be a quasi-compact open.
Then there exists a finitely generated $k$-subalgebra $A' \subset A$
and a quasi-compact open $V' \subset X_{A'}$
such that $V = V'_A$.
\end{lemma}

\begin{proof}
We remark that if $X$ is also quasi-separated this follows from
Limits, Lemma \ref{limits-lemma-descend-opens}. Let
$U_1, \ldots, U_n$ be finitely many affine opens of $X$
such that $V \subset \bigcup U_{i, A}$. Say $U_i = \Spec(R_i)$.
Since $V$ is quasi-compact we can find finitely many
$f_{ij} \in R_i \otimes_k A$, $j = 1, \ldots, n_i$
such that $V = \bigcup_i \bigcup_{j = 1, \ldots, n_i} D(f_{ij})$
where $D(f_{ij}) \subset U_{i, A}$ is the corresponding standard
open. (We do not claim that $V \cap U_{i, A}$ is the union
of the $D(f_{ij})$, $j = 1, \ldots, n_i$.)
It is clear that we can find a finitely generated $k$-subalgebra
$A' \subset A$ such that $f_{ij}$ is the image of some
$f_{ij}' \in R_i \otimes_k A'$.
Set $V' = \bigcup D(f_{ij}')$ which is a quasi-compact open of $X_{A'}$.
Denote $\pi : X_A \to X_{A'}$ the canonical morphism.
We have $\pi(V) \subset V'$ as $\pi(D(f_{ij})) \subset D(f_{ij}')$.
If $x \in X_A$ with $\pi(x) \in V'$, then $\pi(x) \in D(f_{ij}')$
for some $i, j$ and we see that $x \in D(f_{ij})$ as $f_{ij}'$
maps to $f_{ij}$. Thus we see that $V = \pi^{-1}(V')$ as desired.
\end{proof}

\end{document}
